\subsection{Lexical Closures}

A lexical closure exists when an inner function keeps access to names from an enclosing function after the enclosing function has finished running.

This sounds strange only if a function frame is imagined as disappearing completely with all of its local state. In normal cases, a function frame does terminate after return. But when an inner function needs some of the outer function's local bindings, CPython preserves those specific bindings inside cell objects. The whole outer frame is not kept alive as an ordinary active call frame. The necessary referenced state is moved into closure storage.

The result is a function object that carries executable code plus references to preserved lexical state.

\subsubsection{The Lifespan Shift: Preserving Enclosing Environments for Out-of-Scope Execution}

Start with a nested function that is called while the outer function is still running.

\begin{verbatim}
def outer():
    message = "D'oh!"

    def inner():
        return message

    return inner()

result = outer()

print(f"result: {result}")
\end{verbatim}

Possible output:

\begin{verbatim}
result: D'oh!
\end{verbatim}

This example does not yet require a long-lived closure. The inner function runs before \texttt{outer()} returns. The enclosing frame still exists when \texttt{inner()} reads \texttt{message}.

Now return the inner function object itself.

\begin{verbatim}
def outer():
    message = "D'oh!"

    def inner():
        return message

    return inner

fn = outer()

print(f"fn:       {fn}")
print(f"type(fn): {type(fn)}")
print(f"fn():     {fn()}")
\end{verbatim}

Possible output:

\begin{verbatim}
fn:       <function outer.<locals>.inner at 0x...>
type(fn): <class 'function'>
fn():     D'oh!
\end{verbatim}

The function \texttt{outer()} has already returned. Its ordinary execution frame is no longer active. But \texttt{fn()} can still read \texttt{message}. That is the lifespan shift.

The local name \texttt{message} did not simply vanish, because the returned inner function still needs it.

The preserved state is not global.

\begin{verbatim}
def outer():
    message = "D'oh!"

    def inner():
        return message

    return inner

fn = outer()

try:
    print(message)
except NameError as err:
    print(f"error type: {type(err)}")
    print(f"message:    {err}")

print(f"fn(): {fn()}")
\end{verbatim}

Possible output:

\begin{verbatim}
error type: <class 'NameError'>
message:    name 'message' is not defined
fn(): D'oh!
\end{verbatim}

The name \texttt{message} is not available as a global name. It is available only through the closure carried by the returned function object.

A closure can preserve different values for different function objects.

\begin{verbatim}
def make_speaker(message):
    def speak():
        return message

    return speak

say_doh = make_speaker("D'oh!")
say_ni = make_speaker("Ni!")
say_answer = make_speaker("The answer is 42.")

print(say_doh())
print(say_ni())
print(say_answer())
\end{verbatim}

Possible output:

\begin{verbatim}
D'oh!
Ni!
The answer is 42.
\end{verbatim}

Each call to \texttt{make\_speaker()} creates a new function call frame. Each returned \texttt{speak} function preserves the \texttt{message} value from its own call.

The function objects are separate.

\begin{verbatim}
def make_speaker(message):
    def speak():
        return message

    return speak

say_doh = make_speaker("D'oh!")
say_ni = make_speaker("Ni!")

print(f"say_doh:       {say_doh}")
print(f"say_ni:        {say_ni}")
print(f"same object:   {say_doh is say_ni}")
print()
print(f"say_doh():     {say_doh()}")
print(f"say_ni():      {say_ni()}")
\end{verbatim}

Possible output:

\begin{verbatim}
say_doh:       <function make_speaker.<locals>.speak at 0x...>
say_ni:        <function make_speaker.<locals>.speak at 0x...>
same object:   False

say_doh():     D'oh!
say_ni():      Ni!
\end{verbatim}

They were created from the same source code, but they are different function objects with different preserved environments.

The returned function object is still an ordinary function object.

\begin{verbatim}
def make_speaker(message):
    def speak():
        return message

    return speak

fn = make_speaker("No soup for you!")

print(f"type(fn): {type(fn)}")
print()

selected_names = [
    "__call__",
    "__name__",
    "__doc__",
    "__defaults__",
    "__code__",
    "__closure__",
]

for name in selected_names:
    print(f"{name:<16} {name in dir(fn)}")

print()
print(f"fn.__name__: {fn.__name__!r}")
print(f"fn():        {fn()}")
\end{verbatim}

Possible output:

\begin{verbatim}
type(fn): <class 'function'>

__call__        True
__name__        True
__doc__         True
__defaults__    True
__code__        True
__closure__     True

fn.__name__: 'speak'
fn():        No soup for you!
\end{verbatim}

The closure does not make it a different kind of callable. It is still a function object. It just has additional stored lexical state.

Closures are useful when a function must be configured once and executed later.

\begin{verbatim}
def make_prefixer(prefix):
    def add_prefix(text):
        return f"{prefix}: {text}"

    return add_prefix

warn = make_prefixer("WARNING")
info = make_prefixer("INFO")

print(warn("Nobody expects the Spanish Inquisition!"))
print(info("It is only a model."))
\end{verbatim}

Possible output:

\begin{verbatim}
WARNING: Nobody expects the Spanish Inquisition!
INFO: It is only a model.
\end{verbatim}

The function \texttt{make\_prefixer()} builds a specialized function. The returned function remembers the selected \texttt{prefix}.

This is not class-based storage. It is not a global variable. It is a function object carrying references to preserved enclosing bindings.

\subsubsection{How CPython Uses \texttt{\_\_closure\_\_} and Cell Objects to Store Lexical State After the Parent Frame Unwinds}

The closure state of a function object is visible through \texttt{\_\_closure\_\_}.

\begin{verbatim}
def make_speaker(message):
    def speak():
        return message

    return speak

fn = make_speaker("D'oh!")

print(f"fn.__closure__: {fn.__closure__}")
print(f"type:           {type(fn.__closure__)}")
\end{verbatim}

Possible output:

\begin{verbatim}
fn.__closure__: (<cell at 0x...: str object at 0x...>,)
type:           <class 'tuple'>
\end{verbatim}

The \texttt{\_\_closure\_\_} attribute is either \texttt{None} or a tuple of cell objects.

Here it is a tuple containing one cell. That cell contains the preserved object referenced by the outer name \texttt{message}.

The tuple can be inspected.

\begin{verbatim}
def make_speaker(message):
    def speak():
        return message

    return speak

fn = make_speaker("D'oh!")
closure_tuple = fn.__closure__

print(f"type(closure_tuple): {type(closure_tuple)}")
print()

selected_names = [
    "count",
    "index",
    "__len__",
    "__getitem__",
    "__contains__",
]

for name in selected_names:
    print(f"{name:<14} {name in dir(closure_tuple)}")

print()
print(f"len(closure_tuple): {len(closure_tuple)}")
print(f"closure_tuple[0]:   {closure_tuple[0]}")
\end{verbatim}

Possible output:

\begin{verbatim}
type(closure_tuple): <class 'tuple'>

count          True
index          True
__len__        True
__getitem__    True
__contains__   True

len(closure_tuple): 1
closure_tuple[0]:   <cell at 0x...: str object at 0x...>
\end{verbatim}

The cell object exposes its stored object through \texttt{cell\_contents}.

\begin{verbatim}
def make_speaker(message):
    def speak():
        return message

    return speak

fn = make_speaker("D'oh!")
cell = fn.__closure__[0]

print(f"cell:       {cell}")
print(f"type(cell): {type(cell)}")
print()

selected_names = [
    "cell_contents",
    "__class__",
    "__eq__",
    "__repr__",
]

for name in selected_names:
    print(f"{name:<16} {name in dir(cell)}")

print()
print(f"cell.cell_contents: {cell.cell_contents!r}")
print(f"type(contents):     {type(cell.cell_contents)}")
\end{verbatim}

Possible output:

\begin{verbatim}
cell:       <cell at 0x...: str object at 0x...>
type(cell): <class 'cell'>

cell_contents    True
__class__        True
__eq__           True
__repr__         True

cell.cell_contents: "D'oh!"
type(contents):     <class 'str'>
\end{verbatim}

The closure cell stores a reference to the object. It does not store the source-code text \texttt{"message"}. It stores runtime state needed by the inner function.

A function without preserved enclosing names has no closure.

\begin{verbatim}
def plain():
    return "D'oh!"

print(f"plain.__closure__: {plain.__closure__}")
print(f"plain():           {plain()}")
\end{verbatim}

Possible output:

\begin{verbatim}
plain.__closure__: None
plain():           D'oh!
\end{verbatim}

There is no enclosing function state to preserve.

A closure can contain several cells.

\begin{verbatim}
def make_card(name, score):
    mood = "calm"

    def card():
        return f"{name:<10} score={score:04d} mood={mood}"

    return card

fn = make_card("Homer", 42)

print(fn())
print()
print(f"fn.__closure__: {fn.__closure__}")

for index, cell in enumerate(fn.__closure__):
    print(f"{index}: {cell.cell_contents!r}")
\end{verbatim}

Possible output:

\begin{verbatim}
Homer      score=0042 mood=calm

fn.__closure__: (<cell at 0x...: str object at 0x...>, <cell at 0x...: str object at 0x...>, <cell at 0x...: int object at 0x...>)
0: 'calm'
1: 'Homer'
2: 42
\end{verbatim}

The cell order corresponds to the function code object's free-variable order, which will be inspected in the next subsubsection.

If the preserved object is mutable, the closure can mutate that object without rebinding the outer name.

\begin{verbatim}
def make_collector():
    items = []

    def add_item(text):
        items.append(text)
        return items

    return add_item

collector = make_collector()

print(collector("D'oh!"))
print(collector("No soup for you!"))
print(collector("The answer is 42."))
\end{verbatim}

Possible output:

\begin{verbatim}
["D'oh!"]
["D'oh!", 'No soup for you!']
["D'oh!", 'No soup for you!', 'The answer is 42.']
\end{verbatim}

The list object is preserved in the closure. Each call mutates the same list object.

Inspect the preserved list.

\begin{verbatim}
def make_collector():
    items = []

    def add_item(text):
        items.append(text)
        return items

    return add_item

collector = make_collector()

cell = collector.__closure__[0]
items = cell.cell_contents

print(f"items before: {items}")
print(f"id(items):    {id(items)}")
print()

collector("D'oh!")
collector("Ni!")

print(f"items after:  {items}")
print(f"id(items):    {id(items)}")
print(f"cell value:   {cell.cell_contents}")
\end{verbatim}

Possible output:

\begin{verbatim}
items before: []
id(items):    2298078960768

items after:  ["D'oh!", 'Ni!']
id(items):    2298078960768
cell value:   ["D'oh!", 'Ni!']
\end{verbatim}

The same list object remains inside the cell. The list contents change.

The list object can be inspected.

\begin{verbatim}
def make_collector():
    items = []

    def add_item(text):
        items.append(text)
        return items

    return add_item

collector = make_collector()
items = collector.__closure__[0].cell_contents

print(f"type(items): {type(items)}")
print()

selected_names = [
    "append",
    "extend",
    "pop",
    "clear",
    "__len__",
    "__getitem__",
    "__setitem__",
]

for name in selected_names:
    print(f"{name:<14} {name in dir(items)}")

print()
items.append("It is only a model.")
print(f"items: {items}")
\end{verbatim}

Possible output:

\begin{verbatim}
type(items): <class 'list'>

append         True
extend         True
pop            True
clear          True
__len__        True
__getitem__    True
__setitem__    True

items: ['It is only a model.']
\end{verbatim}

If the inner function must rebind an outer name, it needs \texttt{nonlocal}.

\begin{verbatim}
def make_counter():
    count = 0

    def next_count():
        nonlocal count

        count = count + 1
        return count

    return next_count

counter = make_counter()

print(counter())
print(counter())
print(counter())
\end{verbatim}

Possible output:

\begin{verbatim}
1
2
3
\end{verbatim}

The integer object is immutable. The statement:

\begin{verbatim}
count = count + 1
\end{verbatim}

does not mutate the integer object. It creates or obtains a new integer object and updates the closure cell so that \texttt{count} refers to the new object.

Inspect the cell before and after calls.

\begin{verbatim}
def make_counter():
    count = 0

    def next_count():
        nonlocal count

        count = count + 1
        return count

    return next_count

counter = make_counter()
cell = counter.__closure__[0]

print(f"before: {cell.cell_contents}")

counter()
counter()

print(f"after:  {cell.cell_contents}")
\end{verbatim}

Possible output:

\begin{verbatim}
before: 0
after:  2
\end{verbatim}

The closure cell, not a normal local frame slot, is where the preserved binding now lives.

A compact closure structure is:

\begin{verbatim}
outer function call:
    creates local variable message
    creates inner function object
    inner function refers to message
    message is stored in a cell
    inner function receives reference to that cell
    outer function returns inner function

later:
    inner function reads message through the cell
\end{verbatim}

\subsubsection{Inspecting Free Variables via \texttt{co\_freevars}, \texttt{co\_cellvars}, and Closure Cells}

Code objects record closure-related names.

From the outer function's point of view, a local variable that must be preserved for an inner function is a cell variable.

From the inner function's point of view, a variable coming from an enclosing scope is a free variable.

\begin{verbatim}
def outer():
    message = "D'oh!"

    def inner():
        return message

    return inner

fn = outer()

print(f"outer co_cellvars: {outer.__code__.co_cellvars}")
print(f"inner co_freevars: {fn.__code__.co_freevars}")
\end{verbatim}

Possible output:

\begin{verbatim}
outer co_cellvars: ('message',)
inner co_freevars: ('message',)
\end{verbatim}

The same lexical name appears from two perspectives:

\begin{verbatim}
outer:
    message is a cell variable

inner:
    message is a free variable
\end{verbatim}

The inner function's closure tuple stores the cells for its free variables.

\begin{verbatim}
def outer():
    message = "D'oh!"

    def inner():
        return message

    return inner

fn = outer()

free_names = fn.__code__.co_freevars
closure_cells = fn.__closure__

for name, cell in zip(free_names, closure_cells):
    print(f"{name:<10} -> {cell.cell_contents!r}")
\end{verbatim}

Possible output:

\begin{verbatim}
message    -> "D'oh!"
\end{verbatim}

The \texttt{zip()} call pairs each free-variable name with its corresponding closure cell.

For several free variables:

\begin{verbatim}
def make_card(name, score):
    mood = "calm"

    def card():
        return f"{name:<10} score={score:04d} mood={mood}"

    return card

fn = make_card("Homer", 42)

print(f"freevars: {fn.__code__.co_freevars}")
print()

for name, cell in zip(fn.__code__.co_freevars, fn.__closure__):
    print(f"{name:<8} -> {cell.cell_contents!r}")
\end{verbatim}

Possible output:

\begin{verbatim}
freevars: ('mood', 'name', 'score')

mood     -> 'calm'
name     -> 'Homer'
score    -> 42
\end{verbatim}

The exact ordering is the ordering recorded by the code object. Code should not guess it manually. When inspecting, pair \texttt{co\_freevars} with \texttt{\_\_closure\_\_}.

The code object itself can be inspected.

\begin{verbatim}
def make_card(name, score):
    mood = "calm"

    def card():
        return f"{name:<10} score={score:04d} mood={mood}"

    return card

fn = make_card("Homer", 42)
code_obj = fn.__code__

print(f"type(code_obj): {type(code_obj)}")
print()

selected_names = [
    "co_freevars",
    "co_cellvars",
    "co_varnames",
    "co_names",
    "co_consts",
    "co_code",
]

for name in selected_names:
    print(f"{name:<14} {name in dir(code_obj)}")

print()
print(f"co_freevars: {code_obj.co_freevars}")
print(f"co_cellvars: {code_obj.co_cellvars}")
print(f"co_varnames: {code_obj.co_varnames}")
\end{verbatim}

Possible output:

\begin{verbatim}
type(code_obj): <class 'code'>

co_freevars    True
co_cellvars    True
co_varnames    True
co_names       True
co_consts      True
co_code        True

co_freevars: ('mood', 'name', 'score')
co_cellvars: ()
co_varnames: ()
\end{verbatim}

The inner function has no ordinary local variables in this example. Its useful names come from closure cells.

The outer function's code object has cell variables.

\begin{verbatim}
def make_card(name, score):
    mood = "calm"

    def card():
        return f"{name:<10} score={score:04d} mood={mood}"

    return card

outer_code = make_card.__code__

print(f"outer co_varnames: {outer_code.co_varnames}")
print(f"outer co_cellvars: {outer_code.co_cellvars}")
print(f"outer co_freevars: {outer_code.co_freevars}")
\end{verbatim}

Possible output:

\begin{verbatim}
outer co_varnames: ('name', 'score', 'card')
outer co_cellvars: ('name', 'score', 'mood')
outer co_freevars: ()
\end{verbatim}

The outer function has ordinary local names, but some of them also become cell variables because an inner function needs them.

A nested function can have ordinary local variables and free variables at the same time.

\begin{verbatim}
def make_builder(prefix):
    def build(text):
        result = f"{prefix}: {text}"
        return result

    return build

fn = make_builder("INFO")

print(fn("It is only a model."))
print()
print(f"co_varnames: {fn.__code__.co_varnames}")
print(f"co_freevars: {fn.__code__.co_freevars}")
print(f"closure:     {fn.__closure__}")
\end{verbatim}

Possible output:

\begin{verbatim}
INFO: It is only a model.

co_varnames: ('text', 'result')
co_freevars: ('prefix',)
closure:     (<cell at 0x...: str object at 0x...>,)
\end{verbatim}

Here:

\begin{itemize}
    \item \texttt{text} is a parameter local to \texttt{build()};
    \item \texttt{result} is a local variable inside \texttt{build()};
    \item \texttt{prefix} is a free variable coming from the enclosing function;
    \item the actual \texttt{prefix} object is stored in a closure cell.
\end{itemize}

The bytecode can show closure reads.

\begin{verbatim}
import dis

def make_builder(prefix):
    def build(text):
        result = f"{prefix}: {text}"
        return result

    return build

fn = make_builder("INFO")

dis.dis(fn)
\end{verbatim}

Possible output excerpt:

\begin{verbatim}
... LOAD_DEREF               2 (prefix)
... LOAD_FAST                0 (text)
... STORE_FAST               1 (result)
... LOAD_FAST                1 (result)
... RETURN_VALUE
\end{verbatim}

The exact bytecode display varies by Python version. The important distinction is that closure variables are accessed differently from ordinary local variables. A local variable may use a fast-local access path. A closure variable is loaded through a cell.

A closure can be summarized as three connected objects:

\begin{verbatim}
function object
    -> code object
        -> co_freevars names
    -> __closure__ tuple
        -> cell objects
            -> preserved runtime objects
\end{verbatim}

A final inspection example prints the complete mapping.

\begin{verbatim}
def make_counter(start):
    count = start

    def next_count(step=1):
        nonlocal count

        count = count + step
        return count

    return next_count

counter = make_counter(40)

print(counter())
print(counter())
print()

print(f"function object: {counter}")
print(f"code object:     {counter.__code__}")
print(f"freevars:        {counter.__code__.co_freevars}")
print(f"defaults:        {counter.__defaults__}")
print()

for name, cell in zip(counter.__code__.co_freevars, counter.__closure__):
    print(f"{name:<8} -> {cell.cell_contents!r}")
\end{verbatim}

Possible output:

\begin{verbatim}
41
42

function object: <function make_counter.<locals>.next_count at 0x...>
code object:     <code object next_count at 0x..., file "...", line ...>
freevars:        ('count',)
defaults:        (1,)

count    -> 42
\end{verbatim}

The default value \texttt{step=1} is stored in the function object's defaults. The preserved state \texttt{count} is stored through the closure cell. These are different storage mechanisms.

The final rule is:

\begin{quote}
A closure is not a copy of an outer function call. It is a function object whose code refers to preserved cell variables from its lexical enclosing scope.
\end{quote}